\documentclass[12pt]{article}
\usepackage[T1]{fontenc}
\usepackage[utf8]{inputenc}
\usepackage{lmodern}
\usepackage{textcomp}
\usepackage{enumitem}
\usepackage{geometry}
\geometry{margin=1in}
\begin{document}
\begin{center}
Answer Key - Philosophy - Graduate Level Assessment 17 \
\small (Answers are ordered exactly as in the question paper.)
\end{center}

\section*{Multiple Choice}
\begin{enumerate}[label=\arabic*.]
\item \textbf{Answer:} D
\\ \textit{Options:} A) straw man; B) ad populum; C) post hoc fallacy; D) secundum quid; E) ad crumenam
\\ \textit{Note:} The term "secundum quid" refers to an informal fallacy that occurs when a general rule is applied too broadly without sufficient evidence, which aligns with the concept of hasty generalization in logical reasoning.
\item \textbf{Answer:} A
\\ \textit{Options:} A) a good will; B) self-interest; C) justice; D) right action; E) good consequences
\\ \textit{Note:} According to Kant, a good will is intrinsically good because it acts out of a sense of duty and moral obligation, making it the only thing that can be considered good without any qualifications or conditions.
\item \textbf{Answer:} A
\\ \textit{Options:} A) Irrelevant to the truth of the claims; B) Only relevant if the person is a public figure; C) Only valid if the person is conscious of their hypocrisy; D) A sign that the person is untrustworthy; E) Direct evidence of the person's lying tendencies
\\ \textit{Note:} The truth of a claim is determined by its correspondence to reality, independent of the individual's consistency in adhering to their own claims, thus making accusations of hypocrisy irrelevant to the claim's veracity.
\item \textbf{Answer:} D
\\ \textit{Options:} A) \textasciitilde{}Mis; B) Msi\textasciitilde{}; C) M\textasciitilde{}ism; D) \textasciitilde{}Mmis; E) M\textasciitilde{}mis
\\ \textit{Note:} The correct answer \textasciitilde{}Mmis accurately represents the negation of the proposition that Marco moves from Spain to Italy, as it uses the predicate Mxyz to denote the action of moving, where the negation symbol (\textasciitilde{}) indicates that this action does not occur.
\item \textbf{Answer:} A
\\ \textit{Options:} A) Ali; B) Abu Bakr; C) Khadijah; D) Uthman; E) Abu Talib
\\ \textit{Note:} Ali is considered by Shia Muslims to be the first person to recognize Muhammad's prophethood and openly support him, thus playing a crucial role in the early Islamic community's acceptance of Muhammad as a prophet.
\end{enumerate}

\section*{True / False}
\begin{enumerate}[label=\arabic*.]
\setcounter{enumi}{5}
\item \textbf{Answer:} True
\\ \textit{Options:} A) True; B) False
\\ \textit{Note:} Kant defines a hypothetical imperative as a command that applies conditionally, where an action is deemed necessary to achieve a specific end, thus affirming that it represents an action as a means to that desired outcome.
\item \textbf{Answer:} False
\\ \textit{Options:} A) True; B) False
\\ \textit{Note:} Carruther's position actually acknowledges the potential for conscious experience in animals, rejecting the simplistic view that their mental states are solely driven by instinct.
\item \textbf{Answer:} False
\\ \textit{Options:} A) True; B) False
\\ \textit{Note:} The title Dalai Lama actually translates to 'Ocean of Wisdom' in Tibetan, not 'Ocean of Stillness.'
\item \textbf{Answer:} False
\\ \textit{Options:} A) True; B) False
\\ \textit{Note:} The statement is false because the Fatihah is recited in each of the five daily prayers, totaling a minimum of 17 recitations, depending on the number of units (rak'ahs) performed in each prayer.
\item \textbf{Answer:} False
\\ \textit{Options:} A) True; B) False
\\ \textit{Note:} Cicero argues that expedience does not equate to moral rightness, as he emphasizes the importance of virtue and ethical principles that transcend mere practical considerations.
\end{enumerate}

\section*{Long Form Response}
\begin{enumerate}[label=\arabic*.]
\setcounter{enumi}{10}
\item \textbf{Answer:} The laudatory personality fallacy is the tendency to judge an individual's actions based on an idealized perception of their character traits, leading to the assumption that virtuous qualities necessarily correlate with moral behavior. This fallacy has significant implications in moral philosophy, as it can distort ethical evaluations by prioritizing perceived character over the context and consequences of specific actions. Consequently, individuals may be unjustly excused or praised based on their character, rather than the morality of their actions, undermining a more objective ethical analysis. This bias can lead to a failure to hold individuals accountable for harmful actions, as their laudable traits overshadow the moral evaluation of their behavior. Ultimately, recognizing this fallacy is crucial for developing a more nuanced understanding of moral responsibility that transcends simplistic character judgments.
\\ \textit{Note:} The answer correctly identifies the laudatory personality fallacy as a cognitive bias where idealized character traits skew moral judgments, emphasizing the need to consider actions independently of perceived virtues to avoid ethical distortions.
\item \textbf{Answer:} Clarence Darrow's perspective on the meaning of life centers around the notion that life holds intrinsic significance for the majority of individuals, regardless of their circumstances. He believed that meaning arises from personal experiences, relationships, and the pursuit of knowledge, rather than from an abstract or universal truth. This view implies that the significance of life is subjective and can be cultivated through individual engagement with the world. For Darrow, the recognition of life's meaningfulness encourages a compassionate approach to human existence, emphasizing the importance of empathy and understanding in a diverse society. Thus, he posits that embracing our subjective experiences allows for a richer, more fulfilling life for all.
\\ \textit{Note:} The answer correctly captures Darrow's belief that the meaning of life is subjective and derived from personal experiences and relationships, emphasizing that significance is not predetermined but cultivated through individual engagement with one's circumstances.
\end{enumerate}

\vfill
\noindent\textit{End of Answer Key}
\end{document}